\chapter*{KATA PENGANTAR}
\phantomsection
\addcontentsline{toc}{chapter}{KATA PENGANTAR}

% Pedoman 4.3 no. 5: maksimal 2 halaman; terima kasih cukup kepada 5 pihak berikut.
% Ganti teks dalam {[...]} dengan nama lengkap dan gelar. Kurung kurawal wajib selama
% teks masih diawali "[", karena "\item [..." dibaca LaTeX sebagai label nomor.
Puji syukur penulis panjatkan kepada Tuhan Yang Maha Esa karena atas rahmat-Nya penulis dapat menyelesaikan proposal skripsi yang berjudul ``Analisis QoS Jaringan Berbasis \textit{Crowdsourcing} pada Aplikasi Presensi Menggunakan K-Means''. Proposal ini disusun sebagai salah satu syarat untuk mengikuti ujian proposal skripsi pada Program Studi Informatika, Fakultas Teknik dan Informatika, Institut Teknologi dan Bisnis Widya Gama Lumajang.

Penulis mengucapkan terima kasih kepada:
\begin{enumerate}[topsep=0pt,itemsep=0pt]
    \item {[Nama Rektor]}, selaku Rektor Institut Teknologi dan Bisnis Widya Gama Lumajang.
    \item {[Nama Dekan]}, selaku Dekan Fakultas Teknik dan Informatika.
    \item {[Nama Ketua Program Studi]}, selaku Ketua Program Studi Informatika.
    \item {[Nama Dosen Pembimbing]}, selaku Dosen Pembimbing.
    \item Orang tua dan keluarga penulis.
\end{enumerate}

Penulis menyadari bahwa proposal ini masih memiliki kekurangan. Oleh karena itu, kritik dan saran yang membangun sangat penulis harapkan.

\vspace{1cm}
\begin{flushright}
Lumajang, [tanggal] 2026\\[2cm]
Penulis
\end{flushright}
